Kurs:Algebraische Kurven (Osnabrück 2025-2026)/Arbeitsblatt 22/latex
Kurs:Algebraische Kurven (Osnabrück 2025-2026)/Arbeitsblatt 22/latex

\setcounter{section}{22}


  \zwischenueberschrift{Übungsaufgaben}


\inputaufgabe

{}

{


Es sei $R$ der
\definitionsverweis {lokale Ring}{}{}
zum Überkreuzungspunkt des dreidimensionalen Achsenkreuzes. Bestimme dessen
\definitionsverweis {Einbettungsdimension}{}{.}


}

{} {}


\inputaufgabe

{}

{


Man gebe ein Beispiel für eine Kurve


\mavergleichskette

{\vergleichskette

{C
}

{ \subseteq }{ { {\mathbb A}_{  K }^{  n   } }
}

{  }{
}

{    }{
}

{   }{
}

}
{}{}{}
derart, dass es auf ihr Punkte


\mavergleichskette

{\vergleichskette

{ P_1,P_2,P_3
}

{ \in }{ C
}

{  }{
}

{    }{
}

{   }{
}

}
{}{}{}
gibt, deren
\definitionsverweis {Einbettungsdimensionen}{}{}
gleich $1,2,3$ sind.


}

{} {}


\inputaufgabe

{}

{


Es sei


\mavergleichskette

{\vergleichskette

{H(X)
}

{ \in }{ K[X]
}

{  }{
}

{    }{
}

{   }{
}

}
{}{}{,}
sei


\mavergleichskette

{\vergleichskette

{F
}

{ = }{ Y-H
}

{  }{
}

{    }{
}

{   }{
}

}
{}{}{}
und


\mavergleichskette

{\vergleichskette

{C
}

{ = }{ V(F)
}

{ \subseteq }{ {\mathbb A}^{2}_{ K }
}

{    }{
}

{   }{
}

}
{}{}{}
der
\definitionsverweis {Graph}{}{}
von $H$, aufgefasst als
\definitionsverweis {ebene algebraische Kurve}{}{.}
Es sei


\mavergleichskettedisp

{\vergleichskette

{P
}

{ =} { (a,b)
}

{ =} { (a, H(a))
}

{ } {
}

{ } {
}

}
{}{}{}
ein Punkt des Graphen.
\aufzaehlungzweiabc{Zeige, dass die
\definitionsverweis {Multiplizität}{}{}
von $C$ in $P$ gleich $1$ ist.
}{Zeige, dass die
\definitionsverweis {Tangente}{}{}
in $P$ an $C$ mit der üblichen Tangente an einen Graphen im Punkt $a$ übereinstimmt.
}


}

{} {}


\inputaufgabe

{}

{


Es sei $K$ ein
\definitionsverweis {Körper}{}{}
und seien


\mavergleichskette

{\vergleichskette

{ F_1 , \ldots , F_m
}

{ \in }{ K[X_1 , \ldots , X_\ell ]
}

{  }{
}

{    }{
}

{   }{
}

}
{}{}{}
und


\mavergleichskette

{\vergleichskette

{G_1 , \ldots , G_n
}

{ \in }{ K[X_1 , \ldots , X_m ]
}

{  }{
}

{    }{
}

{   }{
}

}
{}{}{}
Polynome, die zu den
\definitionsverweis {polynomialen Abbildungen}{}{}


\mathdisp {{ {\mathbb A}_{ K }^{  \ell  } }  \stackrel{F}{\longrightarrow} { {\mathbb A}_{  K }^{  m   } } \stackrel{G}{\longrightarrow} { {\mathbb A}_{  K }^{  n   } }} {  }


Anlass geben. Es seien
\mathkor {} {J (F)_P} {und} {J (G)_Q} {}
die durch
\definitionsverweis {formales partielles Ableiten}{}{}
definierten
\definitionsverweis {Jacobi-Matrizen}{}{.}
Beweise die formale Kettenregel


\mavergleichskettedisp

{\vergleichskette

{ J ( G \circ F)_P
}

{ =} { J(G)_{F(P)} \circ J(F)_P
}

{ } {
}

{ } {
}

{ } {
}

}
{}{}{.}


}

{} {}


\inputaufgabegibtloesung

{}

{


\aufzaehlungzweiabc{Zeige, dass formales
\definitionsverweis {partielles Ableiten}{}{}
auf dem Polynomring

\mathl{K[X_1 , \ldots , X_n]}{} bezüglich einer Variablen und
\definitionsverweis {Dehomogenisieren}{}{}
bezüglich einer anderen Variablen vertauschbar sind.
}{Zeige, dass dies nicht für die gleiche Variable stimmt.
}


}

{} {}


\inputaufgabegibtloesung

{}

{


Es sei


\mavergleichskette

{\vergleichskette

{ H
}

{ \in }{ K[X_1 , \ldots , X_n]
}

{  }{
}

{    }{
}

{   }{
}

}
{}{}{}
ein
\zusatzklammer {in der
\definitionsverweis {Standardgraduierung}{}{}} {} {}
\definitionsverweis {homogenes Polynom}{}{}
vom Grad $e$. Zeige die Beziehung


\mavergleichskettedisp

{\vergleichskette

{ e H
}

{ =} { X_1 { \frac{   \partial   H   }{  \partial   X_1   } } + \cdots + X_n { \frac{   \partial   H   }{  \partial   X_n   } }
}

{ } {
}

{ } {
}

{ } {
}

}
{}{}{.}


}

{} {}


\inputaufgabe

{}

{


Betrachte die durch


\mavergleichskette

{\vergleichskette

{ y
}

{ = }{ 2x^4+3x^2-x+1
}

{  }{
}

{    }{
}

{   }{
}

}
{}{}{}
gegebene Kurve mit dem Punkt


\mavergleichskette

{\vergleichskette

{ P
}

{ = }{ (1,5)
}

{  }{
}

{    }{
}

{   }{
}

}
{}{}{.}
Finde eine Koordinatentransformation derart, dass $P$ zum Punkt $(0,0)$ wird und die Tangente an $P$ zur $x$-Achse.


}

{} {}


\inputaufgabe

{}

{


Es sei $K$ ein
\definitionsverweis {Körper}{}{}
und


\mavergleichskette

{\vergleichskette

{F
}

{ \in }{ K[X,Y]
}

{  }{
}

{    }{
}

{   }{
}

}
{}{}{}
ein nichtkonstantes
\definitionsverweis {Polynom}{}{}
mit einfachen
\definitionsverweis {Primfaktoren}{}{}
und mit zugehöriger
\definitionsverweis {ebener Kurve}{}{}


\mavergleichskette

{\vergleichskette

{C
}

{ = }{ V(F)
}

{  }{
}

{    }{
}

{   }{
}

}
{}{}{.}
Zeige, dass $C$ nur endlich viele
\definitionsverweis {singuläre Punkte}{}{}
besitzt.


}

{} {}


\inputaufgabegibtloesung

{}

{


Beweise
Lemma 22.12.


}

{} {}


\bild{ \begin{center}

\includegraphics[width=5.5cm]{ \bildeinlesungsvg {Cercle tangente rayon} {svg} }

\end{center}

\bildtext {} }


\bildlizenz { Cercle tangente rayon.svg } {} {} {Commons} {} {}


\inputaufgabegibtloesung

{}

{


Zeige, dass der
\definitionsverweis {Einheitskreis}{}{}
über einem
\definitionsverweis {Körper}{}{}
der
\definitionsverweis {Charakteristik}{}{}
$\neq 2$
\definitionsverweis {glatt}{}{}
ist und bestimme für jeden Punkt die Gleichung der
\definitionsverweis {Tangente}{}{.}


}

{} {}


\inputaufgabe

{}

{


Es sei $K$ ein
\definitionsverweis {Körper}{}{.}
\aufzaehlungzweiabc{Zeige, dass der
\definitionsverweis {Graph}{}{}
eines Polynoms


\mavergleichskette

{\vergleichskette

{F
}

{ \in }{ K[X]
}

{  }{
}

{    }{
}

{   }{
}

}
{}{}{}
eine
\definitionsverweis {glatte}{}{}
\definitionsverweis {algebraische Kurve}{}{}
ist.
}{Es seien


\mavergleichskette

{\vergleichskette

{F,G
}

{ \in }{ K[X]
}

{  }{
}

{    }{
}

{   }{
}

}
{}{}{}
Polynome ohne gemeinsame Nullstelle. Zeige, dass der Graph der
\definitionsverweis {rationalen Funktion}{}{}
$F/G$ ebenfalls eine glatte algebraische Kurve ist.
}


}

{} {}


\inputaufgabegibtloesung

{}

{


Bestimme die
\definitionsverweis {singulären Punkte}{}{}
der
\definitionsverweis {ebenen algebraischen Kurve}{}{}


\mavergleichskettedisp

{\vergleichskette

{ V \left(  -2X^3+3X^2Y-Y+ \frac{2}{3} \sqrt{\frac{1}{3} }  \right)
}

{ \subseteq} { {\mathbb A}^{2}_{ {\mathbb C}}
}

{ } {
}

{ } {
}

{ } {
}

}
{}{}{.}


}

{} {}


\inputaufgabe

{}

{


Bestimme für die in
Beispiel 8.5
berechnete Trajektorie die Koordinaten der Punkte, wo die Kurve singulär ist.


}

{} {}


\inputaufgabegibtloesung

{}

{


Bestimme die Primfaktorzerlegung des Polynoms


\mavergleichskettedisp

{\vergleichskette

{ X^3+XY^2
}

{ \in} { {\mathbb C}[X,Y]
}

{ } {
}

{ } {
}

{ } {
}

}
{}{}{}
und bestimme die Singularitäten der zugehörigen affinen Kurve samt ihren Multiplizitäten und Tangenten.


}

{} {}


\inputaufgabegibtloesung

{}

{


Bestimme die Multiplizität und die Tangenten im Nullpunkt

\mathl{(0,0)}{} der ebenen algebraischen Kurve


\mavergleichskettedisp

{\vergleichskette

{C
}

{ =} { V \left(  Y^4+X^3+3XY^2+2X^2Y  \right)
}

{ \subseteq} { {\mathbb A}^{2}_{ {\mathbb C}}
}

{ } {
}

{ } {
}

}
{}{}{.}


}

{} {}


\inputaufgabegibtloesung

{}

{


Bestimme über die partiellen Ableitungen für das durch das Polynom


\mathdisp {V^3+U^2V-2UV+2U^2-4U-2V} {  }


gegebene Nullstellengebilde einen singulären Punkt. Führe eine Koordinatentransformation durch, die diesen Punkt in den Nullpunkt überführt. Bestimme die Multiplizität und die Tangenten in diesem Punkt.


}

{} {}


\bild{ \begin{center}

\includegraphics[width=5.5cm]{ \bildeinlesungsvg {Cardioid} {svg} }

\end{center}

\bildtext {} }


\bildlizenz { Cardioid.svg } {} {D.328} {Commons} {CC-by-sa 3.0} {}


\inputaufgabe

{}

{


Bestimme die Singularitäten
\zusatzklammer {mit Multiplizitäten und Tangenten} {} {}
der durch


\mathdisp {V \left(  \left(  X^2+Y^2  \right)^2- 2X \left(  X^2+Y^2  \right) -Y^2  \right)} {  }


gegebenen

\definitionswortenp{Kardioide}{.}


}

{} {}


\inputaufgabegibtloesung

{}

{


Betrachte die beiden reellen Kurven


\mathdisp {V \left(  X^5-X^3+2XY+7Y^2-9  \right)} {  }


im Punkt

\mathl{(1,1)}{} und


\mathdisp {V \left(  X^4+Y^4-3X^2Y^2 +5X+7Y  \right)} {  }


im Nullpunkt. Sind diese beiden Kurven lokal in den angegebenen Punkten zueinander
\definitionsverweis {diffeomorph}{}{?}


}

{} {}


  \zwischenueberschrift{Aufgaben zum Abgeben}


\inputaufgabe

{3}

{


Es sei $K$ ein
\definitionsverweis {Körper}{}{}
der
\definitionsverweis {Charakteristik}{}{}


\mavergleichskette

{\vergleichskette

{p
}

{ \geq }{ 0
}

{  }{
}

{    }{
}

{   }{
}

}
{}{}{.}
Man charakterisiere die Polynome


\mavergleichskette

{\vergleichskette

{F
}

{ \in }{ K[X,Y]
}

{  }{
}

{    }{
}

{   }{
}

}
{}{}{}
mit der Eigenschaft, dass
\aufzaehlungdrei{die erste
\definitionsverweis {partielle Ableitung}{}{,}
}{die zweite partielle Ableitung,
}{beide partiellen Ableitungen
}
$0$ sind.


}

{} {}


\inputaufgabe

{4}

{


Bestimme für die
\definitionsverweis {Kurve}{}{}


\mathl{V \left(  X^3+Y^3-3XY+1  \right)}{} die
\definitionsverweis {singulären Punkte}{}{}
über $\R$ und über ${\mathbb C}$. Man gebe jeweils die
\definitionsverweis {Multiplizität}{}{}
und die
\definitionsverweis {Tangenten}{}{}
an.


}

{} {}


\inputaufgabe

{3}

{


Es sei $K$ ein
\definitionsverweis {algebraisch abgeschlossener Körper}{}{}
und seien


\mavergleichskette

{\vergleichskette

{ G,H
}

{ \in }{ K[X,Y]
}

{  }{
}

{    }{
}

{   }{
}

}
{}{}{}
Polynome mit


\mavergleichskette

{\vergleichskette

{ G(P)
}

{ = }{ H(P)
}

{ = }{ 0
}

{    }{
}

{   }{
}

}
{}{}{}
für einen bestimmten Punkt


\mavergleichskette

{\vergleichskette

{ P
}

{ \in }{ {\mathbb A}^{2}_{K}
}

{  }{
}

{    }{
}

{   }{
}

}
{}{}{.}
Es sei


\mavergleichskette

{\vergleichskette

{ F
}

{ = }{ GH
}

{  }{
}

{    }{
}

{   }{
}

}
{}{}{.}
Zeige, dass jede
\definitionsverweis {Tangente}{}{}
von $G$ in $P$ und jede Tangente von $H$ in $P$ auch eine Tangente von $F$ in $P$ ist.


}

{} {}


\inputaufgabe

{6}

{


Es sei $K$ ein
\definitionsverweis {algebraisch abgeschlossener Körper}{}{.}
Betrachte die
\definitionsverweis {Kurve}{}{}


\mavergleichskettedisp

{\vergleichskette

{ C
}

{ =} { V \left(  x^3+5x^2y-6xy^2-x^2-xy+4y^2  \right)
}

{ } {
}

{ } {
}

{ } {
}

}
{}{}{.}
\aufzaehlungdreiabc{Bestimme die
\definitionsverweis {Tangenten}{}{}
im Nullpunkt.
}{Zeige, dass


\mavergleichskettedisp

{\vergleichskette

{ P
}

{ =} { (1,2)
}

{ } {
}

{ } {
}

{ } {
}

}
{}{}{}
ein Punkt der Kurve ist, und berechne die Tangente(n) von $C$ in $P$ über die Ableitung.
}{Führe eine Variablentransformation derart durch, dass $P$ in den neuen Variablen der Nullpunkt ist, und bestimme die Tangente(n) in $P$ aus der transformierten Kurvengleichung.
}


}

{} {}


\inputaufgabe

{4}

{


Bestimme für die
\definitionsverweis {algebraische Kurve}{}{}


\mavergleichskettedisp

{\vergleichskette

{C
}

{ =} { V \left(  9y^4+10x^2y^2+x^4-12y^3-12x^2y+4y^2  \right)
}

{ } {
}

{ } {
}

{ } {
}

}
{}{}{}
die
\definitionsverweis {Singularitäten}{}{}
sowie deren
\definitionsverweis {Multiplizitäten}{}{}
und
\definitionsverweis {Tangenten}{}{.}


}

{\zusatzklammer {Vergleiche dazu
Beispiel 8.5.} {} {}} {}
